\newpage
\phantomsection
\label{prf:q-is-field}

\begin{remark*}[Return]
\hyperref[thm:q-is-field]{Return to Theorem}
\end{remark*}

\begin{theorem*}[$\mathbb{Q}$ Is a Field]
With the operations of
\hyperref[def:rational-operations]{Definition~\ref*{def:rational-operations}},
the system
\[
(\mathbb{Q},+,\cdot,0,1)
\]
is a field.
\end{theorem*}

\begin{proof}
\textbf{Professional Standard Proof.}~\\
By the integral domain theorem, $(\mathbb{Q},+,\cdot,0,1)$ is already an
integral domain. It remains only to verify that every nonzero element has a
multiplicative inverse.

Let $[(a,b)]\in\mathbb{Q}$ with $[(a,b)]\neq 0$. By the zero class
criterion, $a\neq 0$. Since $a\neq 0$ and $b\neq 0$, the pair $(b,a)\in W$,
so $[(b,a)]\in\mathbb{Q}$.

Compute:
\[
[(a,b)]\cdot[(b,a)]
= [(ab,\;ba)]
= [(ab,\;ab)].
\]
Since $ab\neq 0$, we have $(ab,ab)\sim(1,1)$ (as $ab\cdot 1 = 1\cdot ab$),
so $[(ab,ab)] = 1$.

Therefore $[(a,b)]\cdot[(b,a)] = 1$, which means $[(b,a)]$ is the
multiplicative inverse of $[(a,b)]$.

Since every nonzero element of $\mathbb{Q}$ has a multiplicative inverse,
$(\mathbb{Q},+,\cdot,0,1)$ is a field.
\end{proof}

\begin{proof}
\textbf{Detailed Learning Proof.}~\\

\textbf{Step 1.} Recall what must be shown.

By the field definition, a field is an integral domain in which every
nonzero element has a multiplicative inverse. We have already established
that $\mathbb{Q}$ is an integral domain. It therefore suffices to show:
\[
\forall\, [(a,b)]\in\mathbb{Q},\quad [(a,b)]\neq 0
\;\Longrightarrow\;
\exists\, y\in\mathbb{Q},\; [(a,b)]\cdot y = 1.
\]

\textbf{Step 2.} Let $[(a,b)]\in\mathbb{Q}$ be nonzero.

By the zero class criterion, $[(a,b)] = 0$ iff $a = 0$. Since we assumed
$[(a,b)]\neq 0$, we have $a\neq 0$.

\textbf{Step 3.} Propose the candidate inverse.

Define $y := [(b,a)]$. We must first check this is well-formed: since
$a\neq 0$ (established in Step 2) and $b\neq 0$ (since $(a,b)\in W$),
the pair $(b,a)\in W$, so $[(b,a)]\in\mathbb{Q}$.

\textbf{Step 4.} Verify that $[(a,b)]\cdot[(b,a)] = 1$.

By the multiplication formula:
\[
[(a,b)]\cdot[(b,a)]
= [(a\cdot b,\; b\cdot a)]
= [(ab,\; ab)].
\]
We check $(ab,ab)\sim(1,1)$: by the definition of $\sim$, this requires
$ab\cdot 1 = ab\cdot 1$, which is trivially true.

Also, $ab\neq 0$ since $a\neq 0$ and $b\neq 0$ and $\mathbb{Z}$ has no
zero divisors. So $(ab,ab)\in W$ and the class $[(ab,ab)]$ is well-formed.

Since $(ab,ab)\sim(1,1)$, we have $[(ab,ab)] = [(1,1)] = 1$.

Therefore:
\[
[(a,b)]\cdot[(b,a)] = 1.
\]

\textbf{Step 5.} Conclude.

We have exhibited an element $y = [(b,a)]\in\mathbb{Q}$ such that
$[(a,b)]\cdot y = 1$. Since $[(a,b)]$ was an arbitrary nonzero element
of $\mathbb{Q}$, every nonzero element of $\mathbb{Q}$ has a multiplicative
inverse.

Combined with the integral domain theorem, $(\mathbb{Q},+,\cdot,0,1)$ is
a field.
\end{proof}

\begin{remark*}[Proof structure]
The proof has a single substantive step: showing that for a nonzero
$[(a,b)]$, the ``swapped'' class $[(b,a)]$ is a multiplicative inverse.
The key computation is that $[(ab,ab)] = 1$, which follows from the
trivial equivalence $(ab,ab)\sim(1,1)$. The nonvanishing of $ab$ requires
that both $a\neq 0$ (from the nonzero hypothesis) and $b\neq 0$ (from the
well-formedness of the fraction-pair).

The proof builds on the integral domain theorem, which carries all of the
axiom-checking work. The field proof adds only one condition to that package.
\end{remark*}

\begin{remark*}[Dependencies]~\\
\begin{itemize}
  \item \hyperref[def:field]{Definition~\ref*{def:field}}
  \item \hyperref[thm:q-is-integral-domain]{Theorem~\ref*{thm:q-is-integral-domain}}
        \textnormal{($\mathbb{Q}$ is an integral domain)}
  \item \hyperref[thm:ec-zero-numerator]{Theorem~\ref*{thm:ec-zero-numerator}}
        \textnormal{(zero class criterion: $[(a,b)] = 0 \Leftrightarrow a = 0$)}
  \item \hyperref[def:rational-operations]{Definition~\ref*{def:rational-operations}}
        \textnormal{(multiplication formula)}
  \item \hyperref[def:rational-distinguished]{Definition~\ref*{def:rational-distinguished}}
        \textnormal{($1 = [(1,1)]$)}
  \item $\mathbb{Z}$ has no zero divisors (so $a\neq 0$ and $b\neq 0$ implies $ab\neq 0$)
\end{itemize}
\end{remark*}
